\section*{Chapitre \ref{ch:b2:phylogenetic-trees} --- Arbres phylogénétiques et horloge moléculaire}

\begin{solution}{exo:b2:phylogenetic-trees:1}
\omterm{def:b2:phylogenetic-trees:tree}{Homologie}~: un caractère hérité d'un ancêtre commun --- l'humérus, le
radius et l'ulna de l'aile d'une chauve-souris et de la nageoire d'une
baleine. \omterm{def:b2:phylogenetic-trees:tree}{Analogie}~: un caractère semblable acquis indépendamment --- les
ailes des chauves-souris et des oiseaux en tant que surfaces portantes
(les os sont homologues en tant que membres antérieurs, l'aile en tant
que telle ne l'est pas), ou l'œil en chambre noire de la pieuvre et du
vertébré. \omterm{def:b2:phylogenetic-trees:tree}{Synapomorphie}~: une \omterm{def:b2:phylogenetic-trees:tree}{homologie} dérivée partagée par un \omterm{def:b2:phylogenetic-trees:tree}{clade} et
qui le définit --- les plumes pour les oiseaux, l'œuf amniotique pour les
amniotes.
\end{solution}

\begin{solution}{exo:b2:phylogenetic-trees:2}
Oiseaux~: monophylétique. Reptiles sans les oiseaux~: paraphylétique.
Poissons sans les tétrapodes~: paraphylétique (les dipneustes sont plus
proches de nous que de la truite). Vertébrés volants (chauves-souris,
oiseaux, ptérosaures)~: polyphylétique. Mammifères~: monophylétique.
\omterm{def:b2:unicellular-diversity:prokaryote}{Procaryotes}~: paraphylétique (bactéries et archées, sans les eucaryotes
--- et les archées sont plus proches des eucaryotes que des bactéries).
\end{solution}

\begin{solution}{exo:b2:phylogenetic-trees:3}
Groupe frère des oiseaux~: les crocodiles. Groupe frère des mammifères~:
tous les autres amniotes (tortues, lézards, crocodiles et oiseaux
ensemble). Plus petit \omterm{def:b2:phylogenetic-trees:tree}{clade} contenant les lézards et les oiseaux~:
(lézards,(crocodiles, oiseaux)), les diapsides sans les tortues sur cet
arbre. Intervertir crocodiles et oiseaux à leur nœud ne change rien~: on
peut faire pivoter librement un nœud.
\end{solution}

\begin{solution}{exo:b2:phylogenetic-trees:4}
Non enracinés~: $3$~; $3\times 5\times 7 = 105$~; $3\times 5\times
7\times 9\times 11\times 13\times 15 = \num{2027025}$. Arbres enracinés
pour 4 taxons~: chacun des 3 arbres non enracinés a 5 branches, ce qui
en donne 15.
\end{solution}

\begin{solution}{exo:b2:phylogenetic-trees:5}
$((A,B),(C,D))$~: sites 1, 2, 3, un pas chacun, site 4, deux pas
($\mathrm{A}$ et $\mathrm{G}$ des deux côtés), site 5, un~: 6 pas.
$((A,C),(B,D))$~: $2 + 2 + 2 + 1 + 1 = 8$. $((A,D),(B,C))$~: $2 + 2 + 2 +
2 + 1 = 9$. Le premier arbre est le plus parcimonieux~; sur cet arbre, le
site 4 est homoplasique (le changement $\mathrm{A}\to\mathrm{G}$, ou son
inverse, survient deux fois).
\end{solution}

\begin{solution}{exo:b2:phylogenetic-trees:6}
Réunir $A$ et $B$ à la hauteur 3. Puis $d(AB,C) = 14$, $d(AB,D) = 16$,
$d(C,D) = 10$~: réunir $C$ et $D$ à la hauteur 5. Puis $d(AB,CD) = (14 + 16
+ 14 + 16)/4 = 15$~: racine à la hauteur $7.5$. Arbre $((A,B),(C,D))$.
\end{solution}

\begin{solution}{exo:b2:phylogenetic-trees:7}
$d = -0.75\ln(1 - 4p/3)$~: $0.052$, $0.38$, $1.21$. Pour $p = 0.6$,
l'argument du logarithme vaut $0.2$ et la pente $\mathrm{d}d/\mathrm{d}p
= 1/(1 - 4p/3) = 5$~: une erreur d'échantillonnage de $0.02$ sur $p$ devient
$0.1$ sur $d$, et l'hypothèse de vitesses égales en tout site, fausse
pour tout gène réel, ajoute un biais du même ordre. Le gène est saturé.
\end{solution}

\begin{solution}{exo:b2:phylogenetic-trees:8}
$t = d/2k = 0.02/(2\times 10^{-9}) = 10$ millions d'années. À 500
millions d'années, $d = 2\times 10^{-9}\times 5\times 10^{8} = 1.0$ et
$p = 0.75(1 - e^{-4/3}) = 0.55$~: aux trois quarts du chemin vers la
saturation, là où la correction est raide et la date peu fiable
--- il faut un gène plus lent pour de telles séparations.
\end{solution}

\begin{solution}{exo:b2:phylogenetic-trees:9}
\qty{52}{\%}~: le \omterm{def:b2:phylogenetic-trees:tree}{clade} n'est retrouvé que dans à peine la moitié des
rééchantillonnages~; les données ne disent presque rien à son sujet, et
les autres possibilités réunies sont aussi probables. \qty{99}{\%}~: le
signal est cohérent d'un site à l'autre~; le \omterm{def:b2:phylogenetic-trees:tree}{clade} est fortement soutenu
(compte tenu du modèle et de l'alignement --- les erreurs systématiques
comme l'\omterm{prop:b2:phylogenetic-trees:pitfalls}{attraction des longues branches} ne sont pas détectées par le
bootstrap). Pour le premier, ajouter de la séquence (davantage de gènes),
ajouter des taxons qui brisent les branches, et essayer une autre
méthode.
\end{solution}

\begin{solution}{exo:b2:phylogenetic-trees:10}
Deux lignées qui ont chacune changé en une grande fraction de leurs sites
ont, en tout site où les deux ont changé, une chance sur quatre d'aboutir
à la même base (avec quatre bases à vitesses égales). Si chacune a changé
en \qty{40}{\%} des sites, elles ont toutes deux changé en \qty{16}{\%}
et coïncident par hasard en \qty{4}{\%} de l'ensemble des sites --- ce qui
peut dépasser la fraction des sites porteurs des vraies \omterm{def:b2:phylogenetic-trees:tree}{synapomorphies}
qui relient chacune à ses véritables apparentés, à évolution lente. La
parcimonie compte les coïncidences sans se demander pourquoi elles
existent et réunit les deux lignées. Ajouter des taxons qui se détachent
le long de chaque longue branche la divise en segments plus courts, si
bien que les coïncidences fortuites se répartissent sur plusieurs
branches plus courtes et que les vraies \omterm{def:b2:phylogenetic-trees:tree}{synapomorphies} de chaque segment
deviennent visibles~; un modèle de vraisemblance qui prévoit de
nombreuses coïncidences fortuites sur les longues branches les corrige
directement.
\end{solution}

\begin{solution}{exo:b2:phylogenetic-trees:11}
$d = -0.75\ln(1 - 0.12) = 0.096$~; $t = d/2k = 0.096/(4\times 10^{-8})
= 2.4$ millions d'années. Substitutions~: $0.096\times \num{16000} =
1500$, précision relative $1/\sqrt{1500} = \qty{2.6}{\%}$~: une erreur
statistique de $\pm 0.06$ million d'années. Mais la vitesse est étalonnée
sur d'autres séparations et peut être fausse d'un facteur deux pour cette
lignée (variation des vitesses entre lignées, et différence entre la
vitesse généalogique mesurée sur quelques générations et la vitesse
mesurée sur des millions d'années)~; et la divergence du gène est
antérieure à celle des espèces, puisque la population ancestrale était
polymorphe --- la véritable séparation est donc plus récente que celle du
gène, d'un temps qui dépend de la taille de la population ancestrale. (La
date admise, tirée de nombreux gènes nucléaires et des fossiles, est de 6
à 7 millions d'années~: l'horloge mitochondriale est ici trop rapide.)
\end{solution}

\begin{solution}{exo:b2:phylogenetic-trees:12}
\omterm{prop:b2:phylogenetic-trees:pitfalls}{Tri incomplet des lignées}~: lorsque deux spéciations sont proches dans le
temps, un polymorphisme ancestral peut se répartir de telle sorte que
l'arbre d'un gène groupe les espèces autrement que l'arbre des espèces~;
transfert horizontal~: un gène reçu d'une autre lignée porte l'histoire de
cette lignée~; duplication de gènes~: comparer un paralogue chez une
espèce à l'orthologue chez une autre donne la date de la duplication, non
celle de la spéciation. On obtient un arbre d'espèces en utilisant de
nombreux gènes répartis dans tout le génome et en retenant l'arbre que
soutient la majorité (ou un modèle qui prévoit une fraction connue de
gènes discordants), en choisissant soigneusement les orthologues et en
privilégiant, chez les \omterm{def:b2:unicellular-diversity:prokaryote}{procaryotes}, les gènes ribosomiques et
informationnels, rarement transférés.
\end{solution}

\begin{solution}{pb:b2:phylogenetic-trees:1}
\textbf{1.} Sites 1, 2 et 4 (états partagés par $W$ et $H$ seuls,
dérivés par rapport au chameau)~: \omterm{def:b2:phylogenetic-trees:tree}{synapomorphies} de $\{W,H\}$. Sites 3
et 6~: \omterm{def:b2:phylogenetic-trees:tree}{synapomorphies} de $\{W,H,C\}$. Le site 5 est une autapomorphie de
$W$, un changement dans une seule lignée qui ne groupe rien~; aucun site
n'est homoplasique sur l'arbre moléculaire.
\textbf{2.} Un changement par site~: 6 pas.
\textbf{3.} Les sites 1, 2 et 4 demandent désormais deux changements
chacun, les sites 3, 5 et 6 un seul~: 9 pas.
\textbf{4.} L'arbre baleine--hippopotame, avec trois pas d'avance. Trois
sites sur mille constituent une faible marge, que quelques homoplasies
pourraient renverser~; la renforcer demande davantage de gènes, davantage
de taxons (d'autres ruminants, les pécaris, le cerf) et un bootstrap.
\textbf{5.} Le \omterm{def:b2:phylogenetic-trees:tree}{clade} réapparaît dans cinq rééchantillonnages sur six~:
un soutien modéré, bien au-dessus du hasard mais en deçà des
\qty{95}{\%} qu'on qualifie conventionnellement de soutien fort.
\textbf{6.} Le \omterm{meth:b2:phylogenetic-trees:parsimony}{groupe externe} doit se situer hors du groupe, mais assez
près pour que sa séquence reste alignable et non saturée~; le chameau est
un artiodactyle extérieur au groupe vache--porc--hippopotame--baleine. Un
poisson différerait à des niveaux presque saturés, ses états ne diraient
rien de l'état ancestral au sein des mammifères, et sa longue branche
attirerait la lignée du groupe étudié qui évolue le plus vite.
\textbf{7.} $d = 0.0625$, $0.107$, $0.131$, $0.155$.
\textbf{8.} Réunir $W$ et $H$ à la hauteur $0.031$. Puis $d(WH,C) =
0.107$, $d(WH,P) = 0.131$, $d(C,P) = 0.131$, toutes les distances à $M$
valant $0.155$~: réunir $WH$ et $C$ à $0.054$. Puis $d(WHC,P) = 0.131$~:
réunion à $0.065$. Enfin, réunir $M$ à $0.078$. Arbre $(M,(P,(C,(W,H))))$.
\textbf{9.} Oui~: le même arbre, le même emboîtement.
\textbf{10.} Une baleine ayant changé deux fois plus vite serait plus
éloignée de tous les autres, hippopotame compris~; l'UPGMA, qui place
chaque extrémité à la hauteur zéro et réunit la paire la plus proche,
réunirait d'abord l'hippopotame et la vache et laisserait la baleine à
l'extérieur --- exactement l'arbre des anatomistes, pour une mauvaise
raison.
\textbf{11.} $0.0625\times 1000 = 62.5$ substitutions~; écart-type de
Poisson $\sqrt{62.5} = 7.9$.
\textbf{12.} $7.9/62.5 = \qty{13}{\%}$.
\textbf{13.} $k = d/2t = 0.107/(1.2\times 10^{8}) = 8.9\times 10^{-10}$
par site et par an.
\textbf{14.} $t = 0.0625/(2\times 8.9\times 10^{-10}) = 35$ millions
d'années.
\textbf{15.} Porc~: $0.131/(1.79\times 10^{-9}) = 73$ millions d'années~;
chameau~: $0.155/(1.79\times 10^{-9}) = 87$ millions d'années.
\textbf{16.} $35 \pm 4.5$ millions d'années (\qty{13}{\%}).
\textbf{17.} La vitesse varie comme $1/t_{\text{étal}}$ et la date comme
$t_{\text{étal}}$~: $\pm 5/60 = \qty{8}{\%}$, soit $\pm 3$ millions d'années
--- en combinaison avec l'erreur de Poisson, environ $\pm 5$ millions
d'années.
\textbf{18.} $d = -0.75\ln(1 - 0.6) = 0.69$~: le gène a changé, en fait,
en deux tiers de ses sites, la correction a multiplié
$p$ par $1.5$ et sa pente vaut $2.5$~; le gène est proche de la saturation
pour cette séparation et sa date ne vaut pas grand-chose.
\textbf{19.} L'absence d'un caractère n'est pas un état dérivé partagé~:
les baleines ont perdu entièrement le membre postérieur et son astragale,
si bien que l'on ne peut rien dire de l'astragale qu'elles auraient eu.
Une \omterm{def:b2:phylogenetic-trees:tree}{synapomorphie} est un indice en faveur d'un \omterm{def:b2:phylogenetic-trees:tree}{clade}~; sa perte dans une
lignée n'est pas un indice contre l'appartenance à ce \omterm{def:b2:phylogenetic-trees:tree}{clade}.
\textbf{20.} Les baleines fossiles possèdent l'astragale à double poulie~:
les baleines descendent d'un ancêtre qui la possédait, ce qui les place à
l'intérieur des artiodactyles. La prédiction moléculaire est confirmée
par les roches.
\textbf{21.} Paraphylétique~: un ancêtre avec tous ses descendants sauf
la lignée des baleines.
\textbf{22.} \omterm{prop:b2:phylogenetic-trees:pitfalls}{Tri incomplet des lignées} d'un polymorphisme ancestral
lors de deux spéciations rapprochées (vache, hippopotame et baleine se
sont séparés en quelques millions d'années), ou comparaison de
paralogues. Trancher en séquençant de nombreux gènes et en retenant
l'arbre que soutiennent la plupart d'entre eux ainsi que leur
concaténation~; rechercher des duplications dans chaque famille.
\textbf{23.} $d = -0.75\ln(1 - 0.333) = 0.30$~; $k = d/2t =
0.30/(7\times 10^{7}) = 4.3\times 10^{-9}$ par site et par an, cinq
fois la vitesse du gène nucléaire.
\textbf{24.} L'ADN mitochondrial est répliqué par une polymérase à
correction d'épreuves moins efficace, exposé aux radicaux oxygénés de la
chaîne respiratoire et non réparé par la machinerie nucléaire~; il ne
recombine pas non plus, si bien que les dommages s'accumulent.
\textbf{25.} Arbre $(M,(P,(C,(W,H))))$~: les baleines sont le groupe
frère des hippopotames~; $k = 8.9\times 10^{-10}$ par site et par an~;
divergence baleine--hippopotame $35 \pm 5$ millions d'années.
\end{solution}
